\section{超参调优与实验复现}

\subsection{常见超参数组合}

在 CS224R 的实验中，Chelsea 和团队使用以下组合取得了稳定表现：

\begin{itemize}
    \item \textbf{PPO}：clip $\epsilon=0.1$、batch size=64、learning rate=2e-4、value coeff=0.5、entropy coeff=0.01。
    \item \textbf{SAC}：Q learning rate=3e-4、policy learning rate=3e-4、replay buffer size=10e5、batch size=256、target smoothing coeff=0.005。
    \item \textbf{Replay buffer}：prioritized sampling weight $\alpha=0.6$，importance sampling exponent $\beta$ 从 0.4 线性增长到 1.0。
\end{itemize}

\subsection{复现要点}

对抗性地复现实验需要固定随机种子和 deterministic rollout。建议每次 run 记录

\begin{itemize}
    \item seed、env name、max episode length
    \item optimizer state（Adam momenta）与 learning rate schedule
    \item buffer size、prioritization hyperparameters
\end{itemize}

\begin{warningbox}{种子敏感性}
为防止偶发崩溃，必须在 logging 中包括每条 rollout 的 return、length、advantages 的 std。若 return variance 超出预期，应立即停止训练，排查是否 replay buffer 中充斥着 high-loss transitions。
\end{warningbox}

\subsection{本章小结}

超参数配比与记录日志是实验重现的关键。固定 seed、记录 optimizer state、监控 variance 并设置自动中断阈值，能显著降低 long-tail failure 的出现概率。

